% !TEX program = xelatex
% ============================================================================
% ITB-ISLBM 等參形狀函數之 Newton-Raphson 反映射求解
% 專案: 5.Re10595/Edit8_NewInterpolation
% 對應原始碼: itblbm/isoparametric_precompute.h (552 行, 行號以 2026-06-05 版為準)
% 本文件所有數值證據均出自本專案三輪靜態驗證 (2026-06-05):
%   16-agent 工作流 + 獨立第二意見複驗 + 四路總檢驗(含對抗複核)
% 編譯: xelatex newton_raphson_shape_function.tex  (跑兩次以解析交叉引用)
% ============================================================================
\documentclass[11pt,a4paper]{article}

\usepackage[margin=2.4cm]{geometry}
\usepackage{fontspec}
\usepackage{xeCJK}
\setCJKmainfont{Noto Sans CJK TC}
\setCJKsansfont{Noto Sans CJK TC}
\setCJKmonofont{Noto Sans CJK TC}
\usepackage{amsmath,amssymb,bm}
\usepackage{booktabs,array,longtable}
\usepackage{xcolor}
\usepackage{listings}
\usepackage{enumitem}
\usepackage[colorlinks=true,linkcolor=blue!60!black,citecolor=blue!60!black,urlcolor=blue!60!black]{hyperref}

% ---------- listings 設定 (C/CUDA) ----------
\definecolor{codebg}{rgb}{0.97,0.97,0.97}
\definecolor{codekw}{rgb}{0.00,0.30,0.60}
\definecolor{codecm}{rgb}{0.30,0.50,0.30}
\definecolor{codestr}{rgb}{0.55,0.25,0.10}
\lstdefinestyle{cstyle}{
  language=C,
  basicstyle=\ttfamily\footnotesize,
  keywordstyle=\color{codekw}\bfseries,
  commentstyle=\color{codecm},
  stringstyle=\color{codestr},
  numbers=left,
  numberstyle=\tiny\color{gray},
  numbersep=8pt,
  backgroundcolor=\color{codebg},
  frame=single,
  rulecolor=\color{gray!50},
  breaklines=true,
  showstringspaces=false,
  tabsize=4,
  captionpos=b,
  escapeinside={(*@}{@*)},
  morekeywords={size_t,inline,__device__,__forceinline__,__constant__,uchar}
}
\lstset{style=cstyle}

\newcommand{\code}[1]{\texttt{#1}}
\newcommand{\filewhere}[1]{\textcolor{gray}{\small（#1）}}

\title{\textbf{ITB-ISLBM 等參形狀函數之 Newton--Raphson 反映射求解}\\[4pt]
\large 理論、程式實作、逐段解釋與驗證\\[2pt]
\normalsize 專案：\code{5.Re10595/Edit8\_NewInterpolation}（D3Q19 GILBM，週期丘 $Re=10595$）}
\author{Edit8\_NewInterpolation 技術文件\\
\small 對應原始碼：\code{itblbm/isoparametric\_precompute.h}（行號以 2026-06-05 版為準）}
\date{2026 年 6 月 5 日}

\begin{document}
\maketitle

\begin{abstract}
本文件完整說明本專案 ITB-ISLBM（Isoparametric-Transformation-Based,
Interpolation-Supplemented Lattice Boltzmann Method）串流路徑中，
「等參形狀函數反映射」的 Newton--Raphson 求解流程：
給定半拉格朗日出發點（departure foot），在 $7\times 7$ 個節點構成的
六階張量 Lagrange 等參元素上，求解局部座標 $(r,s)$ 使
$\bigl[\,Y(r,s),\,Z(r,s)\,\bigr]=\bigl[\sum_i \phi_i y_i,\ \sum_i \phi_i z_i\bigr]$
命中該出發點。內容涵蓋：(1) 理論——形狀函數構造、反映射非線性方程組、
阻尼 Newton 法與 $2\times 2$ 解析步；(2) 程式碼——與
\code{isoparametric\_precompute.h} 逐行對應的實作列表；(3) 解釋——每段
程式的數值意義與設計取捨；(4) 實作目的與意義——host 端一次性預計算
與 GPU 查表消費（Scheme A）的效能架構；(5) 驗證摘要——本專案三輪
獨立靜態驗證（含有理數精確比對與對抗複核）之結論；以及參考資料。
\end{abstract}

\tableofcontents
\newpage

% ============================================================================
\section{緒論：實作目的與問題定義}
\label{sec:intro}

\subsection{背景}
本專案以曲線貼體網格（body-fitted curvilinear grid）求解週期丘
（periodic hill, $Re=10595$）流場。座標慣例為：
$j$ = 流向（streamwise，MPI 切割、週期）、$k$ = 壁法向（wall-normal，上下壁）、
$i$ = 展向（spanwise，均勻、週期）。在格子 Boltzmann 方法（LBM）中，
非均勻網格上的串流（streaming）步驟需要插值：分布函數 $f_q$ 在時刻
$t+\mathrm{d}t$ 於節點 $\bm{x}_{jk}$ 的值，取自時刻 $t$ 在「出發點」
\begin{equation}
\bm{x}_d=\bm{x}_{jk}-\bm{e}_q\,\mathrm{d}t
\label{eq:foot}
\end{equation}
的插值重建，此即 He--Luo--Dembo 的 interpolation-supplemented LBM
（ISLBM）\cite{HeLuoDembo1996} 與其曲線座標推廣 \cite{HeDoolen1997,Imamura2005}。

ITB-ISLBM \cite{ITBpaper} 的核心想法是：\textbf{不在計算座標系做插值，
而是把出發點直接放在物理座標系，藉由有限元素式的等參變換
（isoparametric transformation）\cite{Hughes2000} 找到它在參考元素中的
局部座標，再用同一組形狀函數產生插值權重}。如此插值算子對網格扭曲
（skewness、stretching）具有一致的幾何意義。

\subsection{要解的方程}
\label{sec:problem}
取以節點 $(j,k)$ 為中心的 $7\times 7$ 物理節點
$\{(y_{ab},z_{ab})\}_{a,b=0}^{6}$（即 $j_0=j-3,\ k_0=k-3$ 起算的方塊），
等參映射定義為
\begin{equation}
\boxed{\;
Y(r,s)=\sum_{a=0}^{6}\sum_{b=0}^{6}L_a(r)\,L_b(s)\;y_{ab},
\qquad
Z(r,s)=\sum_{a=0}^{6}\sum_{b=0}^{6}L_a(r)\,L_b(s)\;z_{ab},
\;}
\label{eq:isomap}
\end{equation}
其中 $\phi_{ab}(r,s)=L_a(r)L_b(s)$ 為張量積形狀函數
（即摘要中 $\sum_i\phi_i x_i$ 的二維張量形式）。
給定出發點 $(y_d,z_d)$（式 \eqref{eq:foot} 的 $y$--$z$ 分量），
\textbf{反映射問題}為：求 $(r,s)$ 使
\begin{equation}
\bm{F}(r,s)\;=\;
\begin{pmatrix} Y(r,s)-y_d \\ Z(r,s)-z_d \end{pmatrix}
\;=\;\bm{0}.
\label{eq:nonlinear}
\end{equation}
因 $L_a$ 為六階多項式，式 \eqref{eq:nonlinear} 是\textbf{無封閉解的
非線性方程組}，本專案以阻尼 Newton--Raphson 法求解
\cite{Press2007,DennisSchnabel1996}。解得 $(r,s)$ 後，插值權重即為
$w_{ab}=L_a(r)L_b(s)$，並（於近壁處）經 ghost 摺疊後固化成查表
（\S\ref{sec:purpose}）。

\subsection{何時解、解幾次}
反映射只在 \textbf{host 端初始化階段執行一次}
\filewhere{main.cu:843--846 呼叫 \code{ITB\_PrecomputeCoefficientsHost}}：
對 9 個 $(e_y,e_z)$ 速度類別（D3Q19 的 19 個 $q$ 依其 $y$--$z$ 投影壓縮
而成；類別 0 即 $e_y{=}e_z{=}0$ 不需插值故跳過）$\times$
內部節點 $j\in[3,\mathrm{NYD6}{-}4],\ k\in[3,\mathrm{NZ6}{-}4]$
各解一次。以現行參數（$\mathrm{NY}=897$、$\mathrm{NZ}=449$、$jp=64$，
故 $\mathrm{NYD6}=21$、$\mathrm{NZ6}=455$）計，每個 rank 共
$8\times15\times449=53{,}880$ 次 Newton 求解、全 64 rank 共
$3{,}448{,}320$ 次。GPU 在其後每一個時間步都\textbf{只查表、不再解任何方程}。

% ============================================================================
\section{理論}
\label{sec:theory}

\subsection{一維七點 Lagrange 形狀函數}
\label{sec:shape}
取等距節點
\begin{equation}
x_n = n-3,\qquad n=0,\dots,6\ \Longrightarrow\ x_n\in\{-3,-2,-1,0,1,2,3\},
\end{equation}
（元素中心 $x=0$ 即當前格點本身，對應索引 $\mathrm{half}=3$）。
標準 Lagrange 基底與其導數為
\begin{align}
L_a(x) &= \prod_{\substack{b=0\\ b\neq a}}^{6}\frac{x-x_b}{x_a-x_b},
\label{eq:lagrange}\\
L_a'(x) &= \sum_{\substack{m=0\\ m\neq a}}^{6}\frac{1}{x_a-x_m}
\prod_{\substack{b=0\\ b\neq a,\,b\neq m}}^{6}\frac{x-x_b}{x_a-x_b}
\qquad\text{（乘積法則逐項展開）}.
\label{eq:dlagrange}
\end{align}
重要性質（本專案以有理數運算 \code{fractions.Fraction} 驗證為
\textbf{精確 0 誤差}，見 \S\ref{sec:verify}）：
\begin{enumerate}[itemsep=1pt]
\item \textbf{單位分割（partition of unity）}：$\sum_a L_a(x)=1$、
      $\sum_a L_a'(x)=0$，對任意 $x$ 成立 —— 保證插值權重和恆為 1
      （質量守恆的必要條件）。
\item \textbf{Kronecker 性質}：$L_a(x_b)=\delta_{ab}$；特別地
      $L(0)=[0,0,0,1,0,0,0]$ —— Newton 失敗時退回 $(r,s)=(0,0)$
      即自動退化為「自我讀取」（identity read）。
\item \textbf{階數}：對任意次數 $\le 6$ 的多項式，值與導數均精確再生；
      $x^7$ 則不再生（負控制已驗，誤差 $\approx 13.8$ 明確非零）——
      插值多項式階數恰為 6，插值誤差收斂率 $O(h^7)$
      \cite{BerrutTrefethen2004}。
\end{enumerate}

\subsection{映射值與 Jacobian 的張量積組裝}
\label{sec:jacobian}
Newton 法需要 $\bm{F}$ 與其 Jacobian。由式 \eqref{eq:isomap} 直接微分：
\begin{equation}
J=\frac{\partial(Y,Z)}{\partial(r,s)}=
\begin{pmatrix} Y_r & Y_s\\ Z_r & Z_s \end{pmatrix},\qquad
\begin{aligned}
Y_r&=\textstyle\sum_{a,b}L_a'(r)L_b(s)\,y_{ab}, &
Y_s&=\textstyle\sum_{a,b}L_a(r)L_b'(s)\,y_{ab},\\
Z_r&=\textstyle\sum_{a,b}L_a'(r)L_b(s)\,z_{ab}, &
Z_s&=\textstyle\sum_{a,b}L_a(r)L_b'(s)\,z_{ab}.
\end{aligned}
\label{eq:J}
\end{equation}
即張量積權重 $N=L_aL_b$、$N_r=L_a'L_b$、$N_s=L_aL_b'$，與形狀函數
共用同一組 $L/L'$ 求值。注意 $r\leftrightarrow a\leftrightarrow j$（流向）、
$s\leftrightarrow b\leftrightarrow k$（法向）的軸對應在整條鏈
（host 求解 $\to$ 儲存 $\to$ device 消費）中保持一致、無交換
（已驗證；若交換，重建出發點誤差會達 $O(1)$）。

\subsection{阻尼 Newton--Raphson 迭代}
\label{sec:newton}
以 $(r_0,s_0)=(0,0)$（元素中心）為初值。每次迭代：
\begin{enumerate}[itemsep=1pt]
\item 求值 $\bm{R}=(R_y,R_z)^\mathsf{T}=\bigl(Y-y_d,\ Z-z_d\bigr)^\mathsf{T}$ 與
      $J$（式 \eqref{eq:J}）。
\item \textbf{$2\times2$ 解析步（Cramer 法）}：
\begin{equation}
\det J = Y_rZ_s-Y_sZ_r,\qquad
\Delta r=\frac{Z_sR_y-Y_sR_z}{\det J},\qquad
\Delta s=\frac{-Z_rR_y+Y_rR_z}{\det J},
\label{eq:cramer}
\end{equation}
可直接驗證 $J\,(\Delta r,\Delta s)^\mathsf{T}=\bm{R}$，
故 $(\Delta r,\Delta s)^\mathsf{T}=J^{-1}\bm{R}$ \textbf{嚴格成立}
（本專案以有理數代數證明，差恰為 0）。不需任何線性代數函式庫。
\item \textbf{阻尼與回溯線搜尋（backtracking line search）}
      \cite{Press2007,DennisSchnabel1996}：
\begin{equation}
(r,s)\ \leftarrow\ (r,s)-\sigma\,(\Delta r,\Delta s),\qquad
\sigma_0=\begin{cases}0.5,&|\Delta r|+|\Delta s|>1\\ 1,&\text{否則}\end{cases}
\end{equation}
若新殘差未下降則 $\sigma\leftarrow\sigma/2$ 重試，直到殘差不增或
$\sigma\le 0.0625$（每個 Newton 步最多 5 次求值）。觸及地板仍接受該步
（非單調容許）——真正不收斂的情形交由失敗路徑處理。
\item \textbf{收斂判據（OR 邏輯，雙單位制）}：
\begin{equation}
\underbrace{\sigma\bigl(|\Delta r|+|\Delta s|\bigr)<10^{-12}}_{\text{局部座標步長}}
\qquad\text{或}\qquad
\underbrace{|R_y|+|R_z|<10^{-11}}_{\text{物理座標殘差}},
\label{eq:converge}
\end{equation}
迭代上限 12 次。$10^{-11}$ 相對最細近壁格距
$\min\Delta z\approx2.256\times10^{-3}$ 為 $\sim4.4\times10^{-9}$
相對殘差，極嚴格。
\item \textbf{奇異護欄}：$|\det J|<10^{-14}$ 即中止並標記失敗。
      健康貼體網格上 $|\det J|\sim O(\Delta y\,\Delta z)\sim10^{-5}$
      （實機最小值 $1.571\times10^{-5}$），與門檻相差約 9 個數量級，
      僅在網格塌縮/翻轉時觸發。
\end{enumerate}
因初值距真解不到半格（$\mathrm{CFL}=0.5\Rightarrow$ 出發點位移
$\le0.5$ 格），Newton 展現標準\textbf{二次收斂}：實測殘差歷程如
$1.57\to0.78\to1.2\times10^{-2}\to1.2\times10^{-6}$，
生產網格上平均 2.0 次、最多 3 次即收斂。

\subsection{失敗處理}
若 12 次未收斂或觸發奇異護欄，求解器回傳 \code{converged=0}，
呼叫端強制 $(r,s)=(0,0)$：由 Kronecker 性質
（\S\ref{sec:shape} 性質 2），權重退化為中心點自我讀取，
物理上等同「該方向不串流」的保守處置。同時在
\code{ITBLBM\_STRICT\_PRECOMPUTE=1}（現行設定）下，任一點失敗即
\code{MPI\_Abort(72)} —— \textbf{絕不靜默使用錯誤流場}。

\subsection{近壁 ghost 外插與權重摺疊}
\label{sec:fold}
$k$ 方向 7 點 stencil 在近壁（$k\in\{3,4,5\}$ 與頂部鏡像）會伸出
實體網格至多 3 層 ghost。處理分兩階段、用\textbf{同一組外插係數}保持自洽：
\begin{enumerate}[itemsep=1pt]
\item \textbf{幾何外插}（Newton 求解時）：ghost 節點座標由過物理列
      $\{3,4,5\}$ 的二次 Lagrange 多項式在 ghost 位置求值
      （\code{GHOST\_EXTRAP\_ORDER=2}）。設 ghost 距第一物理列 $d$ 格，
      外插係數為
\begin{equation}
c_0=\tfrac{(d{+}1)(d{+}2)}{2},\qquad c_1=-d(d{+}2),\qquad c_2=\tfrac{d(d{+}1)}{2},
\qquad c_0+c_1+c_2=1\ \ (\forall d).
\label{eq:ghost}
\end{equation}
\item \textbf{權重摺疊}（求解後）：落在 ghost 列上的插值權重 $w$
      依同一組 $c_i$ 重新分配回物理列。\textbf{摺疊等價定理}
      （本專案數值驗證至 $\sim 9\times10^{-16}$）：對任意資料 $D$，
\begin{equation}
\sum_{b}w_b\,D^{\mathrm{eff}}(k_0^{\mathrm{raw}}+b)
\;=\;
\sum_{b}\tilde w_b\,D(k_0^{\mathrm{phys}}+b),
\label{eq:foldthm}
\end{equation}
其中 $D^{\mathrm{eff}}$ 為 ghost 延拓場、$\tilde w$ 為摺疊後權重。
由式 \eqref{eq:ghost} 係數和為 1，摺疊\textbf{精確保持單位分割}
（實機 $\max\bigl|\sum\tilde w-1\bigr|=1.78\times10^{-15}$）。
\end{enumerate}
故近壁格點的 GPU 表只含物理列索引（$k_{\mathrm{idx}}\in[3,\mathrm{NZ6}-4]$），
kernel 無需任何邊界分支。代價是近壁 $\sim$3 層的 $k$ 方向實質精度
由 degree-6 降為 degree-2（局部、刻意的設計取捨）。

% ============================================================================
\section{程式碼實作}
\label{sec:code}
以下列表均為 \code{itblbm/isoparametric\_precompute.h} 的\textbf{實際程式碼}
（2026-06-05 版，行號經核對零漂移）。

\subsection{節點與形狀函數：\code{ITB\_YZNodeHost} / \code{ITB\_YZShapeHost}}
\filewhere{isoparametric\_precompute.h:18--51；對應式 \eqref{eq:lagrange}--\eqref{eq:dlagrange}}
\begin{lstlisting}[caption={七點 Lagrange 形狀函數及其導數（封閉式直接求值，非解方程）}]
static inline double ITB_YZNodeHost(int n)
{
    return (double)(n - ITB_YZ_ORDER / 2);   /* 7/2=3 -> 節點 {-3..3} */
}

static inline void ITB_YZShapeHost(double x,
                                   double L[ITB_YZ_ORDER],
                                   double dL[ITB_YZ_ORDER])
{
    for (int a = 0; a < ITB_YZ_ORDER; a++) {
        const double xa = ITB_YZNodeHost(a);
        double La = 1.0;                      /* (*@式\eqref{eq:lagrange}@*): L_a = prod (x-xb)/(xa-xb) */
        for (int b = 0; b < ITB_YZ_ORDER; b++) {
            if (b == a) continue;
            const double xb = ITB_YZNodeHost(b);
            La *= (x - xb) / (xa - xb);
        }
        L[a] = La;

        double dLa = 0.0;                     /* (*@式\eqref{eq:dlagrange}@*): 乘積法則, 對被微分項 m 求和 */
        for (int m = 0; m < ITB_YZ_ORDER; m++) {
            if (m == a) continue;
            const double xm = ITB_YZNodeHost(m);
            double term = 1.0 / (xa - xm);
            for (int b = 0; b < ITB_YZ_ORDER; b++) {
                if (b == a || b == m) continue;
                const double xb = ITB_YZNodeHost(b);
                term *= (x - xb) / (xa - xb);
            }
            dLa += term;
        }
        dL[a] = dLa;
    }
}
\end{lstlisting}

\subsection{映射與 Jacobian 組裝：\code{ITB\_EvaluateMapHost}}
\filewhere{isoparametric\_precompute.h:87--117；對應式 \eqref{eq:isomap}、\eqref{eq:J}}
\begin{lstlisting}[caption={張量積一次組裝 $Y,Z$ 與四個 Jacobian 分量（每次 Newton 迭代呼叫一次）}]
static inline void ITB_EvaluateMapHost(
    const double *y_h, const double *z_h,
    int j0, int k0, double r, double s,
    double *Y, double *Z,
    double *Yr, double *Ys, double *Zr, double *Zs)
{
    double Lr[ITB_YZ_ORDER], dLr[ITB_YZ_ORDER];
    double Ls[ITB_YZ_ORDER], dLs[ITB_YZ_ORDER];
    ITB_YZShapeHost(r, Lr, dLr);              /* r 方向: 值 + 導數 */
    ITB_YZShapeHost(s, Ls, dLs);              /* s 方向: 值 + 導數 */

    *Y = 0.0; *Z = 0.0;
    *Yr = 0.0; *Ys = 0.0; *Zr = 0.0; *Zs = 0.0;
    for (int a = 0; a < ITB_YZ_ORDER; a++) {
        const int gj = j0 + a;                /* r <-> a <-> j(流向) */
        for (int b = 0; b < ITB_YZ_ORDER; b++) {
            const int gk = k0 + b;            /* s <-> b <-> k(法向) */
            const double yv = ITB_GeomEffHost(y_h, gj, gk);  /* 近壁 ghost 幾何 */
            const double zv = ITB_GeomEffHost(z_h, gj, gk);  /* 由二次外插補出  */
            const double N  = Lr[a]  * Ls[b];   /* phi_ab            */
            const double Nr = dLr[a] * Ls[b];   /* d(phi_ab)/dr       */
            const double Ns = Lr[a]  * dLs[b];  /* d(phi_ab)/ds       */
            *Y  += N  * yv;   *Z  += N  * zv;
            *Yr += Nr * yv;   *Ys += Ns * yv;
            *Zr += Nr * zv;   *Zs += Ns * zv;
        }
    }
}
\end{lstlisting}

\subsection{Newton 求解器：\code{ITB\_NewtonSolveHost}}
\filewhere{isoparametric\_precompute.h:119--180；對應 \S\ref{sec:newton} 全部}
\begin{lstlisting}[caption={阻尼 Newton--Raphson 反映射求解器（完整原文）}]
static inline int ITB_NewtonSolveHost(
    const double *y_h, const double *z_h,
    int j0, int k0, double yd, double zd,
    double *r_out, double *s_out,
    int *iters_out, double *res_out, double *update_out, double *min_det_out)
{
    double r = 0.0, s = 0.0;                  /* 初值 = 元素中心(格點本身) */
    double min_det = 1.0e300;
    double last_update = 0.0;
    double res_norm = 1.0e300;
    int converged = 0;
    int it = 0;

    for (it = 1; it <= 12; it++) {            /* 迭代上限 12 */
        double Y, Z, Yr, Ys, Zr, Zs;
        ITB_EvaluateMapHost(y_h, z_h, j0, k0, r, s, &Y, &Z, &Yr, &Ys, &Zr, &Zs);
        const double Ry = Y - yd;             /* 殘差 F = map - target */
        const double Rz = Z - zd;
        const double det = Yr * Zs - Ys * Zr; /* det J */
        const double abs_det = fabs(det);
        if (abs_det < min_det) min_det = abs_det;
        res_norm = fabs(Ry) + fabs(Rz);
        if (abs_det < 1.0e-14) break;         /* 奇異護欄 -> 標記失敗 */

        double dr = ( Zs * Ry - Ys * Rz) / det;   /* (*@式\eqref{eq:cramer}@*): Cramer */
        double ds = (-Zr * Ry + Yr * Rz) / det;   /* = J^{-1} R     */
        double scale = 1.0;
        last_update = fabs(dr) + fabs(ds);
        if (last_update > 1.0) scale = 0.5;   /* 大步先砍半 */

        double best_r = r - scale * dr;       /* Newton 方向: x <- x - J^{-1}F */
        double best_s = s - scale * ds;
        double best_res = 1.0e300;
        for (int damp = 0; damp < 5; damp++) {    /* 回溯線搜尋 */
            double Yc, Zc, Yrc, Ysc, Zrc, Zsc;
            ITB_EvaluateMapHost(y_h, z_h, j0, k0, best_r, best_s,
                                &Yc, &Zc, &Yrc, &Ysc, &Zrc, &Zsc);
            best_res = fabs(Yc - yd) + fabs(Zc - zd);
            if (best_res <= res_norm || scale <= 0.0625) break;
            scale *= 0.5;
            best_r = r - scale * dr;
            best_s = s - scale * ds;
        }

        r = best_r;
        s = best_s;
        last_update *= scale;                 /* = 實際施加步長 |dr|+|ds| 乘 sigma */
        res_norm = best_res;
        if (last_update < 1.0e-12 || res_norm < 1.0e-11) {   /* (*@式\eqref{eq:converge}@*) 雙判據 */
            converged = 1;
            break;
        }
    }

    *r_out = r;
    *s_out = s;
    *iters_out = it;
    *res_out = res_norm;
    *update_out = last_update;
    *min_det_out = min_det;
    return converged;
}
\end{lstlisting}

\subsection{呼叫端：出發點、求解、權重固化}
\filewhere{isoparametric\_precompute.h:465--510（節錄）}
\begin{lstlisting}[caption={每個 (class, j, k) 解一次；解完立刻把 $(r,s)$ 固化成權重表}]
for (int j = 3; j < NYD6 - 3; j++) {
    for (int k = 3; k < NZ6 - 3; k++) {
        const int j0 = j - ITB_YZ_ORDER / 2;          /* stencil 基底 = j-3, k-3 */
        const int k0 = k - ITB_YZ_ORDER / 2;
        const double yd = y_h[j * NZ6 + k] - ey * dt_val;   /* (*@式\eqref{eq:foot}@*): 出發點 */
        const double zd = z_h[j * NZ6 + k] - ez * dt_val;
        /* ... */
        const int ok = ITB_NewtonSolveHost(y_h, z_h, j0, k0, yd, zd,
                                           &r, &s, &iters, &residual,
                                           &update, &min_det);
        /* ...統計記錄... */
        if (!ok) {                                    /* 失敗 fallback:        */
            c->flags |= ITB_COEFF_NEWTON_FAILED;      /* (0,0) -> Kronecker    */
            r = 0.0;                                  /* -> identity read      */
            s = 0.0;                                  /* (strict 模式仍會 abort)*/
        }

        double wr[ITB_YZ_ORDER], dwr[ITB_YZ_ORDER];
        double ws_raw[ITB_YZ_ORDER], dws[ITB_YZ_ORDER];
        double ws_folded[ITB_YZ_ORDER];
        ITB_YZShapeHost(r, wr, dwr);                  /* w_a = L_a(r)  (j 方向) */
        ITB_YZShapeHost(s, ws_raw, dws);              /* w_b = L_b(s)  (k 方向) */
        for (int a = 0; a < ITB_YZ_ORDER; a++) c->wr[a] = wr[a];
        unsigned char fold_flags = 0;
        ITB_FoldKWeightsHost(k, ws_raw, c->k_idx, ws_folded, &fold_flags);
        c->flags |= fold_flags;                       /* 近壁 ghost 權重摺回物理列 */
        for (int b = 0; b < ITB_YZ_ORDER; b++) c->ws[b] = ws_folded[b];
        /* ... */
    }
}
\end{lstlisting}

\subsection{GPU 消費端：純查表（Scheme A，不重算）}
\filewhere{isoparametric\_streaming.h:42--53（$e_x=0$ 路徑節錄）}
\begin{lstlisting}[caption={device 每時間步只做查表乘加；Newton 與形狀函數一生只在初始化跑一次}]
const ITB_YZCoeff c = itb_yz_coeff_d[ITB_CoeffIndexDevice(yz_id, j, k)];
if (ex == 0.0) {
    double out = 0.0;
    for (int sj = 0; sj < ITB_YZ_ORDER; sj++) {
        const int gj = c.j0 + sj;             /* j 方向: 直接索引(MPI halo) */
        const double wj = c.wr[sj];
        for (int sk = 0; sk < ITB_YZ_ORDER; sk++) {
            const int gk = c.k_idx[sk];       /* k 方向: 摺疊後物理列索引 */
            out += wj * c.ws[sk] *
                   f_post_read[q_off + gj * nface + gk * NX6 + i];
        }
    }
    return out;
}
\end{lstlisting}

% ============================================================================
\section{逐段解釋}
\label{sec:explain}

\subsection{為什麼殘差/Jacobian 的搭配是對的}
程式定義 $\bm{F}=\text{map}-\text{target}$（\code{Ry=Y-yd}），
故 Newton 更新必須是 $\bm{x}\leftarrow\bm{x}-J^{-1}\bm{F}$ ——
程式的 \code{r - scale*dr} 正是此符號方向。式 \eqref{eq:cramer} 的
$2\times2$ Cramer 展開已用有理數代數逐項驗證
$J(\Delta r,\Delta s)^{\mathsf T}-\bm{R}=\bm{0}$（差恰為 0）；
若軸向交換（$r\leftrightarrow s$ 接錯），數值實驗顯示重建出發點誤差
立即升到 $O(1)$，現行程式為正確慣例。

\subsection{阻尼策略的三層保險}
\begin{enumerate}[itemsep=1pt]
\item \textbf{先驗砍半}：首步 $|\Delta r|+|\Delta s|>1$（出發點超過一格）
      即 $\sigma_0=0.5$，避免六階多項式在元素外的劇烈振盪把迭代甩飛。
\item \textbf{回溯線搜尋}：殘差不降即 $\sigma$ 對半，至多 5 次求值——
      這是教科書式 globally convergent Newton 的簡化版
      \cite{Press2007}。
\item \textbf{非單調地板}：$\sigma\le0.0625$ 仍接受該步。看似危險，
      實則安全：真正不收斂者必以 \code{converged=0} 結束，而
      strict 模式會 \code{MPI\_Abort} —— 失敗不可能漏網。
      實機上地板從未觸及（平均 2.0 步收斂）。
\end{enumerate}

\subsection{收斂判據為何用兩種單位}
$|\Delta r|+|\Delta s|<10^{-12}$ 量的是\textbf{局部（元素）座標}，
天然無因次、與格距無關；$|R_y|+|R_z|<10^{-11}$ 量的是\textbf{物理座標}
絕對殘差。OR 邏輯讓兩者互為備援：步長極小但殘差量測受捨入污染時
由前者收斂，反之由後者。理論保證不會「假收斂」：
$\|\Delta\|\ge\|R\|/\|J\|$ 而本網格 $\|J\|\sim$ 格距 $\ll1$，
故步長判據觸發時殘差必然更小（對抗驗證 \S\ref{sec:verify} 已代數排除
false-converge）。

\subsection{初值為何取 $(0,0)$ 而不用更聰明的猜測}
出發點位移 $|\bm{e}|\,\mathrm{d}t$ 在 $\mathrm{CFL}=0.5$ 下不超過半格，
即真解必在 $|r|,|s|\lesssim0.5$ 內（實機統計 $|r|,|s|>1$ 的計數為 0）。
元素中心已是「半徑 0.5 內」的猜測，配合二次收斂 2--3 步即達
$10^{-12}$，再聰明的初值已無利可圖。

\subsection{失敗 fallback 的因果鏈}
$(r,s)=(0,0)\ \xrightarrow{\text{Kronecker}}\ L(0)=e_3
\ \xrightarrow{\text{摺疊}}\ $ 內部節點 $kg=k\ge3$ 走 interior 分支
$\Rightarrow$ 摺疊後唯一非零權重 $=1$ 且 \code{k\_idx} 指回 $k$ 自身
—— 即使在近壁，identity read 也精確指回原格點（已逐案驗證
$k\in\{3,4,5,\dots\}$）。這條鏈的正確性依賴 \S\ref{sec:shape}
性質 2 與摺疊在 interior 分支的恆等性。

\subsection{已知的非數值性瑕疵（誠實記錄）}
三輪驗證唯一確認的問題：\textbf{迭代計數 off-by-one}
\filewhere{isoparametric\_precompute.h:132,175} ——
12 次全跑完未 break 時，C 語言 for 迴圈後置遞增使 \code{it}=13，
\code{*iters\_out} 回報 13（實跑 12 次）。此值\textbf{只}流入
\code{newton\_iter\_sum}/\code{newton\_max\_iter\_used} 兩個診斷計數器
（precompute.h:478,481），\textbf{絕不}流向 $r,s,w_r,w_s,k_{\mathrm{idx}}$；
且該路徑必伴隨 \code{newton\_failed>0}，strict 模式會中止整個 run，
故僅在「已中止 run 的最後一段 log」可見。屬休眠的統計瑕疵，
非數值錯誤。修正建議：\code{*iters\_out = (it > 12 ? 12 : it);}。

% ============================================================================
\section{實作目的與意義}
\label{sec:purpose}

\subsection{Host 一次性預計算 vs.\ kernel 每步重算}
反映射 Newton 解一次約需 2--3 次映射求值（每次 $49$ 節點 $\times$ 6 個量
的張量積），若放在 GPU kernel 每時間步重算，相當於每步每點每方向多付
數百次乘加與分支。本實作把\textbf{全部}非線性求解移到 host 初始化
（每 rank $53{,}880$ 次、瞬間完成），GPU 端時間迴圈內：
\begin{center}
\textbf{0 次 Newton、0 次形狀函數求值、0 個邊界分支} —— 只有
$7{+}49{+}343$ 級的查表乘加。
\end{center}
這延續本專案集的共通預計算哲學（MRT $K$ 矩陣、eta 權重共享，
見專案 \code{CLAUDE.md}「GILBM 效能優化架構」）。

\subsection{Scheme A：存權重而非存 $(r,s)$}
歷史上曾實作 Scheme B（每點只存 16 B 的 $(r,s)$、device 重算形狀函數
與摺疊），後改為現行 \textbf{Scheme A}（commit \code{77eafb8}）：
直接存 \code{ITB\_YZCoeff\{j0, k\_idx[7], wr[7], ws[7], flags\}}，
device 純查表。實測 \textbf{+37\% MLUPS}（達 907 MLUPS/GPU），
代價是係數表增至每 rank $\sim$12.5 MB —— 在 H200 上以記憶體換計算
明確划算。形狀函數的「求解」與「求值」因此徹底離線化。

\subsection{9 類別壓縮}
插值權重只依賴 $(e_y,e_z)$ 而非完整 19 個 $q$（$y,z$ 與 $x$ 無關、
$x$ 方向另以 2 組固定權重 \code{ITB\_WX} 處理），故表壓縮為
9 個類別：表大小與 Newton 求解次數皆省 53\%。
類別映射已對全部 19 個 $q$ 與 \code{GILBM\_e} 三方比對一致。

\subsection{方法學意義}
相對於計算座標插值（GILBM 原路徑），ITB 等參反映射的意義在於：
(1) 插值權重有\textbf{明確的物理空間幾何意義}（出發點真的落在
$L_a(r)L_b(s)$ 的單位分割上），對強拉伸貼體網格不引入度規近似誤差；
(2) 權重的所有幾何複雜度（曲線網格、近壁 ghost、MPI halo、週期接縫）
都被凍結在一張表裡，kernel 與邊界邏輯完全解耦；
(3) 失敗模式安全——不收斂即中止，不會默默退化精度。

% ============================================================================
\section{驗證摘要}
\label{sec:verify}
本文件所有主張均出自 2026-06-05 的三輪獨立靜態驗證
（多代理工作流 $\to$ 獨立第二意見 $\to$ 四路總檢驗含對抗複核），
三輪結論一致：\textbf{正確（correct）}。要點：

\begin{center}
\begin{tabular}{@{}p{0.42\textwidth}p{0.52\textwidth}@{}}
\toprule
\textbf{驗證項} & \textbf{結果（容差 $10^{-12}$）}\\
\midrule
$L_a$、$L_a'$ vs.\ 精確多項式微分 & 有理數比對，誤差恰為 $0$（bit-exact）\\
$\sum L=1$、$\sum L'=0$、Kronecker & 有理數精確成立\\
degree $\le6$ 再生 / degree-7 不再生 & 精確 $0$ / 明確非零（負控制通過）\\
$2\times2$ 步 $=J^{-1}\bm{F}$、更新符號 & 有理數代數證明，差 $=0$\\
合成彎曲映射回收已知 $(r^*,s^*)$ & 最大誤差 $4.1\times10^{-12}$，二次收斂\\
摺疊等價定理（式 \eqref{eq:foldthm}） & $\sim9\times10^{-16}$；單位分割保持\\
ghost 外插係數 vs.\ 二次 Lagrange & 精確相等；$\sum c_i=1\ \forall d$\\
類別映射 / host--device 索引 / 消費域 & 19 個 $q$ 全一致；索引公式字元級相同；消費域 $=$ 填表域\\
\midrule
實機統計（64 GPU 生產網格） &
$3{,}448{,}320/3{,}448{,}320$ 全收斂；平均/最多 $2.0/3$ 次；
$\max$ 殘差 $9.999\times10^{-12}$；$\min|\det J|=1.571\times10^{-5}$；
$\max|\sum\tilde w-1|=1.78\times10^{-15}$\\
\bottomrule
\end{tabular}
\end{center}

對抗複核另外\textbf{排除}了兩個表面疑慮（det 護欄不設 converged、
摺疊和警告用 rank 本地值）——均證明在本網格與本演算法結構下不可達。
唯一確認的瑕疵為 \S4.6 的診斷統計 off-by-one。

\paragraph{限制與注意事項}
(1) 近壁 $\sim$3 層 ghost 帶的 $k$ 方向實質精度為 degree-2；
(2) 7 點半寬恰好吃滿 $\mathrm{Buffer}=3$ 的 MPI halo——升 9 點需先加寬 halo；
(3) \code{ITB\_FillCenterCoeff} 的 clamp 預填在消費域外是「死資料」，
若未來擴大 kernel $(j,k)$ 範圍或讓類別 0 改查表，須先修正其 identity
指向（潛在地雷，已記錄於專案記憶）。

% ============================================================================
\begin{thebibliography}{9}
\bibitem{ITBpaper}
\emph{An Isoparametric Transformation-Based Interpolation-Supplemented
Lattice Boltzmann Method}（ITB-ISLBM 原始論文；本專案實作之基準文獻，
專案內部參考檔名 \code{1.AN\_ISOPARAMETRIC\_TRANSFORMATION-BASED\_INTERPOLATIO.pdf}）.

\bibitem{HeLuoDembo1996}
X.~He, L.-S.~Luo, M.~Dembo,
``Some progress in lattice Boltzmann method. Part~I. Nonuniform mesh grids,''
\emph{Journal of Computational Physics} \textbf{129}, 357--363 (1996).
（ISLBM：插值補充式 LBM 之原始文獻）

\bibitem{HeDoolen1997}
X.~He, G.~Doolen,
``Lattice Boltzmann method on curvilinear coordinates system:
flow around a circular cylinder,''
\emph{Journal of Computational Physics} \textbf{134}, 306--315 (1997).
（曲線座標 ISLBM）

\bibitem{Imamura2005}
T.~Imamura, K.~Suzuki, T.~Nakamura, M.~Yoshida,
``Acceleration of steady-state lattice Boltzmann simulations on
non-uniform mesh using local time step method,''
\emph{Journal of Computational Physics} \textbf{202}, 645--663 (2005).
（GILBM 系譜：廣義座標 LBM）

\bibitem{Hughes2000}
T.~J.~R.~Hughes,
\emph{The Finite Element Method: Linear Static and Dynamic Finite
Element Analysis}, Dover Publications (2000).
（等參元素與形狀函數之標準教材）

\bibitem{Press2007}
W.~H.~Press, S.~A.~Teukolsky, W.~T.~Vetterling, B.~P.~Flannery,
\emph{Numerical Recipes: The Art of Scientific Computing}, 3rd ed.,
Cambridge University Press (2007), \S9.6--9.7.
（含線搜尋之全域收斂 Newton 法）

\bibitem{DennisSchnabel1996}
J.~E.~Dennis~Jr., R.~B.~Schnabel,
\emph{Numerical Methods for Unconstrained Optimization and Nonlinear
Equations}, SIAM Classics in Applied Mathematics (1996).
（阻尼 Newton 與回溯線搜尋之理論基礎）

\bibitem{BerrutTrefethen2004}
J.-P.~Berrut, L.~N.~Trefethen,
``Barycentric Lagrange interpolation,''
\emph{SIAM Review} \textbf{46}(3), 501--517 (2004).
（Lagrange 插值之數值性質）

\bibitem{Frohlich2005}
J.~Fr\"ohlich, C.~P.~Mellen, W.~Rodi, L.~Temmerman, M.~A.~Leschziner,
``Highly resolved large-eddy simulation of separated flow in a channel
with streamwise periodic constrictions,''
\emph{Journal of Fluid Mechanics} \textbf{526}, 19--66 (2005).
（週期丘基準算例與網格幾何）
\end{thebibliography}

\end{document}
